\section{Recurrence, bias selection and empirical laws}\label{sec:laws}
\subsection{Every mixture recurs}
\begin{proof}[Proof of Theorem~\ref{thm:cluster}]
For every rational pair $0<u<v<1$, Theorem~\ref{thm:lower} makes
$\sum_n\Pp_{1/2}(u<T_n<v)=\infty$. Write $h_n$ for their hit probabilities. These events depend only on the disjoint effective collars $C_n$ and are independent. For completeness, independence gives
\[
 \Pp(\text{no hit from }N\text{ through }M)
 =\prod_{n=N}^M(1-h_n)\le\exp\left(-\sum_{n=N}^M h_n\right)\longrightarrow0.
\]
Thus each rational interval is hit infinitely often, almost surely. Intersect these events over their countable basis. On this one probability-one event, every $\lambda\in(0,1)$ is approached by $T_{n_j}$ along an increasing subsequence, by choosing successively smaller rational intervals and a later hit in each. Lemma~\ref{lem:calibration} implies $\mu_{n_j}^\eta\Rightarrow\nu_\lambda$. Every cluster state is DLR and hence belongs to \eqref{eq:DLR}. Finally the cluster set of a sequence in a compact metric space is closed: a diagonal choice of increasingly late points within shrinking distances of converging cluster points proves this directly. It therefore also contains $\nu_0$ and $\nu_1$.
\end{proof}

\subsection{Biased boundaries}
\begin{proof}[Biased part of Theorem~\ref{thm:laws}]
At $p_n=1/2+n^{-2/5}$, Lemmas~\ref{lem:drift} and \ref{lem:transfer} imply, for every fixed $d>0$ and all large $n$,
\[
 \Pp_{p_n}(X_n<a_+-d)
 \le Cn^2e^{-c n^{2/5}}+e^{-c'n^{1/5}}.
\]
Indeed positive $F_n$ of order $n^{3/5}$ has logistic value tending to one, and the transfer error vanishes. This bound is summable.

Fix $p>1/2$. Eventually $p\ge p_n$. Couple exterior words coordinatewise by shared uniforms; ferromagnetic monotonicity makes $X_n$ increasing. Its lower-tail probability at $p$ is bounded by that at $p_n$. Summability implies that $X_n\ge a_+-d$ eventually almost surely: the probability of any failure after size $N$ is at most the tail sum, tending to zero. Intersect over positive rational $d$. For every boundary, $X_n$ is bounded above by its all-plus boundary value, which tends to $a_+$ by the definition of $\mu^+$. Hence $X_n\to a_+$ almost surely. Lemma~\ref{lem:calibration} gives $\mu_n^\eta\Rightarrow\mu^+$. Spin reversal treats $p<1/2$, including $p=0$; $p=1$ is the constant-plus boundary itself.
\end{proof}

\subsection{Fair concentration near the pure phases}
\begin{lemma}[Fair anti-concentration]\label{lem:anticoncentration}
For every nonnegative sequence $r_n=o(\sqrt n)$,
\[
 \Pp_{1/2}(|F_n|\le r_n)\longrightarrow0.
\]
\end{lemma}
\begin{proof}
Fix $\varepsilon>0$ and select $n$ of the north positions in Lemma~\ref{lem:tooth}. Call a position good when $d_v\ge m/2$. A bad position forces $R_{n,v}^++R_{n,v}^->m/2$ by \eqref{eq:gap}. For a sufficiently large fixed even $D=D(\varepsilon)\ge12$, \eqref{eq:rem} makes each expected remainder eventually below $\varepsilon m/16$. Markov's inequality therefore shows that the event $G_D$ of at least $n/2$ good positions satisfies
\[
 \Pp_{1/2}(G_D^c)\le
 \frac4{nm}\sum_v\E(R_{n,v}^++R_{n,v}^-)\le\varepsilon
\]
for large $n$. No independence among these positions is needed.

Clip $F_n$ to $[-r_n-4K,r_n+4K]$, obtaining an increasing function $f$. Whenever $|F_n(\eta)|\le r_n$, the two values of every bit pair stay within this clipping interval; hence $\osc_i f=\osc_i F_n$ there. On that event intersected with $G_D$, Lemma~\ref{lem:response} gives
\[
 \sum_i\osc_i f\ge(n/2)\gamma_Dm/2=n\gamma_Dm/4.
\]
Let $S$ count all $M_n$ exterior plus bits. Under the fair product law, conditioning on the other coordinates gives
\[
 \sum_i\E\osc_i f=4\Cov(f,S)
 \le4\sqrt{\Var(f)\Var(S)}
 \le2(r_n+4K)\sqrt{M_n}\le6(r_n+4K)\sqrt n.
\]
All oscillations outside the specified event are nonnegative. Therefore
\[
 \Pp_{1/2}(|F_n|\le r_n)
 \le\varepsilon+\frac{24(r_n+4K)}{\gamma_Dm\sqrt n}.
\]
Take $n\to\infty$ for this fixed $D$, and then $\varepsilon\downarrow0$. This order of limits is what permits $\gamma_D$ to be small.
\end{proof}

\begin{lemma}[Law of a single random state]\label{lem:law}
Under the fair boundary law,
\[
 \mathcal L(\mu_n^\eta)\Rightarrow
 \tfrac12\delta_{\mu^+}+\tfrac12\delta_{\mu^-}.
\]
\end{lemma}
\begin{proof}
Oddness of $F_n$ and Lemma~\ref{lem:anticoncentration}, with constant sequences $r_n$, put asymptotic probability $1/2$ on each of its two diverging signs. Lemma~\ref{lem:transfer} and $\Pp(H_n^c)\to0$ imply
$\mathcal L(X_n)\Rightarrow\tfrac12\delta_{a_+}+\tfrac12\delta_{a_-}$.
Uniform calibration, Lemma~\ref{lem:calibration}, transfers this to the state space: $\pi(T_n)$ has asymptotic law $\tfrac12\delta_1+\tfrac12\delta_0$, and the distance of $\mu_n^\eta$ from $\nu_{\pi(T_n)}$ vanishes uniformly.
\end{proof}

\begin{proof}[Fair part of Theorem~\ref{thm:laws}]
We give the bounded independent strong law needed for the empirical assertion. If $(Y_n)$ are independent, centred, and bounded by a common constant, then
$\Var(\sum_{n\le N}Y_n)\le C N$. At $N=k^2$, Chebyshev bounds the probability of $|k^{-2}\sum_{n\le k^2}Y_n|>d$ by $C/(d^2k^2)$. The summable-tail argument gives convergence to zero along squares, on intersecting over positive rational $d$. Between $k^2$ and $(k+1)^2$, boundedness changes the normalized sums by $O(1/k)$. Thus $N^{-1}\sum_{n\le N}Y_n\to0$ almost surely.

Apply this to a rational polynomial $P$ of finitely many fixed cylinder probabilities, evaluated at $\mu_n^\eta$, and subtract its expectation. Once its cylinder supports lie in $B_n$, the evaluation depends only on $C_n$, so these bounded variables are independent after finitely many terms. By Lemma~\ref{lem:law}, their expected empirical mean tends to $[P(\mu^+)+P(\mu^-)]/2$. Intersect the resulting probability-one events over the countable algebra of rational cylinder polynomials. This algebra determines weak convergence of laws on $\mathcal P(\Omega)$: it contains the constants, separates states, and its real uniform closure is all continuous functions by the Stone--Weierstrass theorem. Equivalently, uniform continuity first reduces to finitely many cylinder coordinates and polynomial approximation on their compact cube finishes the reduction. We obtain the empirical convergence asserted in Theorem~\ref{thm:laws}.

For the mean of $X_n$, spin reversal gives $\E_{1/2}X_n=1/2$ exactly at every size. The variables $X_n$ are independent and bounded, so the same strong law directly gives their ordinary average $1/2$. This last calculation in fact needs no low-temperature condition. For the spin expectation rather than the plus-probability, the average is zero.
\end{proof}
